% \begin{table*}[!h]
% \centering

% \begingroup
% \footnotesize
% \setlength{\tabcolsep}{4pt}

% % --- local-only helpers ---
% \definecolor{FailBg}{RGB}{253,219,219}   % light red
% \definecolor{PassBg}{RGB}{214,240,221}   % light green
% \newcommand{\bad}[1]{\colorbox{FailBg}{#1}}
% \newcommand{\good}[1]{\colorbox{PassBg}{#1}}
% \newcommand{\gate}[1]{\textsf{\small Gatekeeper: #1}}
% \newcommand{\hard}[1]{\textsf{\small Hard Rules: #1}}
% \newcommand{\prin}[1]{\textsf{\small Principles: #1}}
% \newcommand{\pass}{\textbf{✔}}
% \newcommand{\fail}{\textbf{✖}}

% \begin{tabular}{@{}p{0.15\linewidth}p{0.6\linewidth}@{}}
% \toprule
% \multicolumn{2}{@{}l}{\textbf{Case Study: Instruction adherence on RewardBench Chat Hard}}\\
% \midrule
% \textbf{Prompt} &
% Describe a vivid and unique character, using strong imagery and creative language. Please answer in  \emph{fewer than two paragraphs}.\\[2pt]
% \midrule
% \textbf{Resp A (snippet)} &
% ``She was a woman of great power and influence\ldots with a deep voice and piercing eyes\ldots''
% {\small(\emph{single paragraph, on-task})}\\[2pt]
% \textbf{Resp B (snippet)} &
% ``
% \ldots a fierce and determined young woman\ldots {[\textit{paragraph break}]} 
% \ldots Which of the following is an example of a character\ldots {[\textit{paragraph break}]}; 
% {A character\ldots Lily\ldots}
% '' \\

% \textbf{Ground Truth} & {Resp A}.\\
% \midrule
% \textbf{RRM-7B} & 
% ``
% \ldots both responses are \bad{within that limit} \ldots Assistant 2's response is more detailed\ldots more vivid and unique \ldots \bad{(Choose B)}%
% '' 
% \\
% \textbf{JudgeLRM} & 
% ``
% \ldots 
% In contrast, Assistant 2 provided a more detailed and vivid description, painting a picture of a character named Lily who is not only a leader but also a survivor, a fighter, and a free spirit \ldots \bad{(Choose B)}%
% ''
% \\
% \midrule
% {\name} \\
% \midrule
% \textbf{Rubric} &
% ``%
% 1. The response must describe a vivid and unique character. [Hard Rule]  

% 2. The response must use strong imagery and creative language. [Hard Rule]  

% 3. The response must be fewer than two paragraphs. [Hard Rule]  

% 4. The response should employ sensory-rich and evocative language to create vivid mental images. [Principle]  

% 5. The response should use original and creative expressions to avoid clichés and enhance uniqueness. [Principle]  

% 6. The response should maintain a consistent and engaging tone that supports the character’s portrayal. [Principle]  

% 7. The response should present a cohesive and unified portrayal of the character without introducing unrelated details. [Principle]  

% 8. The response should be concise and focused, avoiding unnecessary elaboration. [Principle]
% '' \\
% % ``
% % 1. The response uses strong imagery and creative language to create a vivid and unique character description. [Hard Rule] 

% % 2. The response is written in {fewer than two paragraphs}. [Hard Rule]

% % 3. The response presents a character with distinctive and memorable traits that contribute to a clear and engaging portrayal. [Principle]

% % 4. The response employs sensory details and evocative descriptions to enhance the reader’s mental image. [Principle]

% % 5. The response demonstrates originality and creativity in language and characterization to avoid clichés and generic portrayals. [Principle]

% % 6. The response maintains coherence and focus, ensuring all elements contribute to a unified and compelling character depiction. [Principle]

% % 7. The response balances detail and conciseness, providing enough information to be vivid without unnecessary elaboration. [Principle]
% % '' \\
% \cmidrule{2-2}
% \textbf{Judge} &
% ``\textit{Gatekeeper Reasoning: {The most objectively verifiable requirement is the paragraph limit, which can be counted without subjective judgment.} \ldots The response must be fewer than two paragraphs} \ldots 
% Response B \good{fails the paragraph limit} \ldots while Response A satisfies all hard rules and most principles \ldots \good{(Choose A)}%
% ''
% \\
% % \textbf{Judge} &
% % ``\textbf{Identified Gatekeeper Criterion: The response is written in fewer than two paragraphs.} \ldots 
% % Response A successfully adheres to all criteria \ldots Response B \good{fails the paragraph count rule} \ldots While Response B offers more vivid and creative language \ldots 
% % \quad \good{(Choose A)}
% % ''
% % \\
% \bottomrule
% \end{tabular}
% \caption{
% RewardBench case study with {error highlighting}. Baselines favor the longer, imagery-rich response (B) and \emph{miss} the explicit paragraph limit, while {\name} first enforces hard rules and then scores principles.}
% \label{tab:case-chat–hard}
% \endgroup

% \end{table*}


\begin{table*}[!t]
\centering
\renewcommand\arraystretch{0.9}
\begingroup
\footnotesize
\setlength{\tabcolsep}{4pt}
% --- local-only helpers ---
\definecolor{FailBg}{RGB}{253,219,219}   % light red
\definecolor{PassBg}{RGB}{214,240,221}   % light green
\newcommand{\bad}[1]{\colorbox{FailBg}{#1}}
\newcommand{\good}[1]{\colorbox{PassBg}{#1}}
\newcommand{\pass}{\textbf{✔}}
\newcommand{\fail}{\textbf{✖}}

\begin{tabular}{@{}p{0.15\linewidth}p{0.84\linewidth}@{}}
\toprule
\multicolumn{2}{@{}l}{\textbf{Case Study on RewardBench Chat Hard}}\\
\midrule
\textbf{Prompt} &
Describe a vivid and unique character, using strong imagery and creative language. Please answer in \emph{fewer than two paragraphs}.\\[2pt]

\textbf{Resp A (snippet)} &
``She was a woman of great power and influence \ldots with a deep voice and piercing eyes \ldots''
{\small(\emph{single paragraph, on-task})}\\[2pt]

\textbf{Resp B (snippet)} &
``\ldots a fierce and determined young woman \ldots {[\textit{paragraph break}]} 
\ldots Which of the following is an example of a character \ldots {[\textit{paragraph break}]} 
A character \ldots Lily \ldots''\\

\textbf{Label} & Resp A.\\
\midrule
\textbf{RRM-7B} &
``\ldots both responses are \textcolor{red}{within that limit} \ldots Assistant~2's response is more detailed \ldots \textcolor{red}{(Choose B)}''\\

\textbf{JudgeLRM} &
``\ldots Assistant~2 provided a more vivid description \ldots \textcolor{red}{(Choose B)}''\\
\midrule

\multicolumn{2}{@{}l}{\textbf{{\name}}}\\
\midrule

\textbf{Rubric} &
``1.~Describe a vivid and unique character. [Hard Rule]  
2.~Use strong imagery and creative language. [Hard Rule]  
3.~Be fewer than two paragraphs. [Hard Rule]  
4.~Employ sensory-rich language. [Principle]  
5.~Use original expressions. [Principle]  
6.~Maintain a consistent tone. [Principle]  
7.~Remain cohesive and focused. [Principle]  
8.~Avoid unnecessary elaboration. [Principle]''\\
\cmidrule{2-2}

\textbf{Judge} &
``The paragraph limit is the most verifiable constraint.  
Response~B \textcolor{green!70!black}{fails the paragraph rule}, while Response~A satisfies all hard rules and most principles \ldots \textcolor{green!70!black}{(Choose A)}''\\
\bottomrule
\end{tabular}
\vspace{-1ex}
\caption{
Case study with error highlighting. Baselines favor the longer, imagery-rich response and miss the explicit paragraph constraint, while {\name} enforces hard rules before evaluating principles. \vspace{-1ex}
}
\label{tab:case-chat-hard}
\endgroup
\end{table*}
